%!TEX encoding = UTF-8 Unicode

%% Thesis main document

% Remove 'coverdepname' parameter to remove department name from cover pages
% Add 'final' parameter to remove "DOCUMENTO PROVISÓRIO" from cover page
\documentclass[12pt,a4paper,twoside,openany]{memoir}

%%% This file is a template for the cover of your
%%%  thesis.
%%% Lines like this one, starting with %%% , are
%%%  instructions. Search them to know where and
%%%  how to edit this file. After editing the
%%%  document you can delete them.
%%% Items marked with OPTIONAL can be removed.
%%% Most of the lines you will have to change
%%%  already have examples to help you understand.
%%% The complete list of options and commands can
%%%  be found in uaThesisTemplate.sty

%%% Select your area of study and department from
%%%  the following lists and replace them in the
%%%  \usepackage declaration ([area,department]).
%%%
%%% | area        |   | department |
%%% ---------------   --------------
%%% | arts        |   | dao        |
%%% | sciences    |   | dbio       |
%%% | education   |   | dcspt      |
%%% | economy     |   | deca       |
%%% | engineering |   | degeit     |
%%% | humanities  |   | de         |
%%% | health      |   | deti       |
%%% |             |   | demac      |
%%% |             |   | decivil    |
%%% |             |   | dem        |
%%% |             |   | dfis        |
%%% |             |   | dlc        |
%%% |             |   | dmat       |
%%% |             |   | dq         |
%%% |             |   | dcm        |
%%% |             |   | dgeo        |
%%% ---------------   --------------
%% Import the UA thesis template
\usepackage[engineering,deca]{uaThesisTemplate}

%% Optional packages for utf8 encoding and T1 fonts
\usepackage[utf8]{inputenc}

\usepackage{palatino}       % Loads Palatino as the serif (roman) font
\usepackage{mathpazo}  % math & text
\renewcommand\familydefault{\rmdefault} % Use serif font (Palatino) as default
\usepackage[T1]{fontenc}

%% We will use english and portuguese
\usepackage[english,portuguese]{babel}

\usepackage{xspace}% used by \sigla

%%%%%%%%%%%%%%%%%%%%%%%%%%%%%%%%%%%%%%%%%%%%%%%%%%
%%       DEFINE THESIS COVER FIELDS
%%                  START
%%%%%%%%%%%%%%%%%%%%%%%%%%%%%%%%%%%%%%%%%%%%%%%%%%

%%% Edit thesis year
%% Thesis Year
\ThesisYear{\the\year}


%%% Uncomment one of the following options:
%\DocumentTypeThesis
%\DocumentTypeThesisProposal
\DocumentTypeDissertation

%%% Uncomment if you are a student of a MAP Doctoral program 
%%% http://www.map.edu.pt/
%%% | area                | identifier |
%%% | Biology of Plants   | bioplant   |
%%% | Computer Science    | i          |
%%% | Physics             | fis        |
%%% | Telecommunications  | tele       |

%\MAPProgram{i}

%%% Edit author's name
%% Author's full name
\Author{Bernardo}{Figueiredo}

%%% Edit thesis title in Portuguese
%% Thesis Mandatory Title in Portuguese
% TODO: draft translation, revisit wording once scope is final
\TitlePT{Visualização Generalizada de Sistemas de Digital Twin}

%%% OPTIONAL: you can write your thesis title in
%%%  other languages. You can repeat the following
%%%  line as many times as you want, or remove it
%%%  if you only want the portuguese title.
%% Optional Titles in additional languages.
% TODO: working title, revisit once scope is final
\TitleEN{Generalized Viewing of Digital Twin Systems}

%%% Edit your supervisor's information.
%%%
%%% Use this command for male supervisor
%%% \Supervisor{name}{category}
%%%
%%% Or this command for female supervisor
%%% \Supervisora{name}{category}
%%%
%%% name   : Supervisor's name
%%% categoy: One of the following list
%%%
%%% | category                  |
%%% -----------------------------
%%% | catedrático/a             |
%%% | catedrático/a convidado/a |
%%% | associado/a c/ agregação  |
%%% | associado/a               |
%%% | auxiliar                  |
%%% | auxiliar convidado        |
%%% -----------------------------
%% Supervisor information
\Supervisor{Diogo Nuno Pereira Gomes}{associado}
% Needed because the student (DeCA) and the supervisor (DETI) are not from the
% same department: without it, the template would list the supervisor under DeCA.
\SupervisorDepartment{Departamento de Eletr\'{o}nica, Telecomunica\c{c}\~{o}es e Inform\'{a}tica}

%%% OPTIONAL: you can add a second supervisor.
%%%
%%% Use this command for male supervisor
%%% \SupervisorMember{name}{category}
%%% 
%%% Or this command for female supervisor
%%% \SupervisoraMember{name}{category}
%%%  {[preposition: (at, in, no, da)] department}
%%%  {[preposition: (at, in, do, na)] school}
%%%  {[preposition: (and, e)] }
%%%
%%% name      : See above.
%%% category  : See above.
%%% department: Can be a department, faculty or
%%%             other; Choose the preposition
%%%             accordingly.
%%% school    : University/School name; Choose the
%%%             preposition accordingly.
%%% last      : Write {e } (Portuguese) or {and } (English)
%%%             for the last enumerated Supervisor; Or
%%%             nothing {} for all other supervisors 
%% Optional SecondSupervisor information
% Co-supervisor. Added directly to the list because \SupervisorMember always
% prints "Professor <category>", which does not fit a researcher.
\listadd{\ciSupervisorMember}{, e do Doutor João Almeida, Investigador do Instituto de Telecomunicações}

%%% Edit your degree information.
%%%
%%% \Degree{title}{name}
%%%
%%% title: Degree title (Mestre or Doutor)
%%% name : Degree name
%% Degree title and name
\Degree{Mestre}{Desenvolvimento de Jogos Digitais}

%%% OPTIONAL: edit grant support texts where
%%%  appropriate. Maximum two.
%%%
%%% \GrantText{grant1}{grant2}
%% Optional Grant support texts
% No grant support to declare.
%\GrantText{...}{...}

%%% Edit the Jury President information.
%% First Member will be the President of the Jury
% TODO: jury not assigned yet, placeholder to fill in once known
\JuriMember{Prof. X}{professor associado da Universidade de Aveiro}

%%% Repeat the above declaration for additional
%%%  members of the jury as many times as necessary.
%% Additional members of the Jury
\JuriMember{Prof. Doutor Diogo Nuno Pereira Gomes}{professor associado da Universidade de Aveiro}
% TODO: add remaining jury members once known

%%% OPTIONAL: Add an acknowledgement text
%% Optional Acknowledgement text
\Acknowledgements{Agradeço à minha família e aos meus amigos pelo apoio ao longo deste trabalho.%
  \\Agradeço também a todos os que me ajudaram pelo caminho e, em particular, às pessoas que dedicaram o seu tempo a testar a aplicação e a partilhar a sua opinião.
}

%%% Edit the keywords and abstract in Portuguese
%%% Don't forget to remove the \blindtext[3]
%%%  declaration.
%%% Don't remove the \selectlanguage statement or you will end with incorrect hyphenation
%%% Don't use commands in the Abstract. Only the \selectlanguage{}, \texttt{}, and \\ commands were configured to be supported. You can configure additional commands in \pdfstringdefDisableCommands
%% Mandatory Keywords and Abstract in Portuguese
\Abstract{palavras-chave}{gémeos digitais, Unreal Engine, visualização 3D, dados geoespaciais, arquitetura de software, infraestrutura rodoviária, código aberto.}
         {resumo}{\selectlanguage{portuguese}As infraestruturas rodoviárias estão cada vez mais instrumentadas com sensores e complementadas por simulações, o que gera grandes volumes de dados heterogéneos que os operadores precisam de interpretar no seu contexto espacial. Os gémeos digitais tornam isso possível através de representações virtuais atualizadas em tempo real. No entanto, as aplicações que os visualizam são, em regra, desenvolvidas como soluções fechadas e específicas, presas a um único tipo de infraestrutura e a uma única plataforma de dados (backend). Esta dissertação, desenvolvida no âmbito do projeto DT4MOB, apresenta um visualizador de gémeos digitais de código aberto, construído sobre o Unreal Engine e o Cesium for Unreal, que é generalizável a diferentes tipos de infraestrutura e independente da plataforma de dados que o alimenta. Todo o conhecimento sobre a plataforma de dados fica confinado a uma camada de serviços, isolado por interfaces reduzidas, e o conhecimento específico de cada tipo de infraestrutura fica concentrado num modelo de entidades assente numa tabela de registo e em dois pontos de extensão opcionais. Deste modo, a renderização, a interação, a câmara e o posicionamento geoespacial são implementados uma única vez. O visualizador suporta cerca de trinta tipos de entidades, distribuídos por cinco casos de uso dos parceiros do projeto (monitorização geotécnica, simulação de incêndios, tráfego, gestão de ativos rodoviários e portagens), sem código específico de cada tipo nos processos genéricos de criação, atualização e renderização. Uma avaliação de desempenho de três estratégias de renderização para um tipo de entidade muito numeroso, a sinalização vertical, mostra que uma abordagem sem agrupamento de instâncias não escala e que o instanciamento por textura tem, na prática, o mesmo custo por frame que um atlas de texturas, ocupando muito menos memória. Um estudo de usabilidade com dez participantes da população-alvo obteve uma pontuação média de 81,1 na System Usability Scale, com os principais problemas concentrados nos controlos de navegação e na visibilidade das entidades, e não na arquitetura.}

%%% OPTIONAL: you can write the keywords and
%%%  abstract in other languages. You can repeat
%%%  the following lines as many times as you want
%%%  or remove them if you only want portuguese.
%%% Don't remove the \selectlanguage statement or you will end with incorrect hyphenation.
%%% Don't use commands in the Abstract. Only the \selectlanguage{}, \texttt{}, and \\ commands were configured to be supported. You can configure additional commands in \pdfstringdefDisableCommands
%%% You can modify the language to any other supported by Latex
%% Optional Keywords and Abstract in additional languages.
\Abstract{keywords}{digital twins, Unreal Engine, 3D visualisation, geospatial data, software architecture, road infrastructure, open source.}
         {abstract}{\selectlanguage{english}Road infrastructure is increasingly instrumented with sensors and complemented by simulations, producing large volumes of heterogeneous data that operators need to see in context. Digital twins make this possible through live virtual representations, but their viewers are typically built as closed, purpose-specific applications, tied to a single infrastructure type and a single data backend. This dissertation, developed within the DT4MOB project, presents an open-source digital twin viewer built on Unreal Engine and Cesium for Unreal that is generalisable across infrastructure types and decoupled from its data backend. Backend knowledge is confined to a service layer behind small interfaces, and infrastructure-specific knowledge to an entity model driven by a registration table and two optional extension points, so that rendering, interaction, the camera, and geospatial placement are written once. The viewer supports roughly thirty entity types across five partner use cases, including geotechnical monitoring, wildfire simulation, traffic, road asset management, and tolling, with no type-specific code in its generic spawn, update, or rendering paths. A benchmark of three rendering strategies for a high-cardinality entity type, road signs, shows that an unbatched approach does not scale and that per-texture instancing is practically on par with a texture atlas in frame time while using far less memory. A user study with ten participants from the intended user population produced a mean System Usability Scale score of 81.1, with the main issues concentrated in the navigation controls and entity visibility, not in the architecture.}


%%% OPTIONAL: Acknowledge the use of AI Tools
\AIUse{I acknowledge the use of Claude (Anthropic, https://claude.ai/) to summarise notes and articles and proofread drafts. I acknowledge the use of Grammarly (Grammarly Inc., https://www.grammarly.com) for grammar and clarity suggestions on the text, each of which was reviewed before being accepted.
}
%Reconheço o uso de [inserir sistema(s) de IA e respetiva hiperligação] para [indicar utilização específica de IA ou outras tarefas]. Reconheço a utilização de [indicar software, código ou plataforma] para [indicar utilização específica de software, código ou platafoma ou para outras tarefas].
%Exemplo 1: Reconheço a utilização do ChatGPT 3.5 (Open AI, https://chat.openai.com) para resumir as notas iniciais e para rever o rascunho final, a utilização do Office365 (Microsoft, https://www.office.com) para a escrita do texto.
%Exemplo 2: Não foram utilizados no presente trabalho quaisquer conteúdos gerados por tecnologias de IA.

%%%%%%%%%%%%%%%%%%%%%%%%%%%%%%%%%%%%%%%%%%%%%%%%%%
%%       DEFINE THESIS COVER FIELDS
%%                  END
%%%%%%%%%%%%%%%%%%%%%%%%%%%%%%%%%%%%%%%%%%%%%%%%%%


%%%%%%%%%%%%%%%%%%%%%%%%%%%%%%%%%%%%%%%%%%%%%%%%%%
%%       BEGIN DOCUMENT PACKAGES
%%%%%%%%%%%%%%%%%%%%%%%%%%%%%%%%%%%%%%%%%%%%%%%%%%
% Enable for for Overleaf support
%\usepackage{lmodern}

\usepackage{ifthen}
\def\useoverleaf{0}  % change to non-zero (for instance, 1) to enable it

\pdfpxdimen=\dimexpr 1in/300\relax % Set DPI to 300 for better quality

\usepackage{morewrites}  % Allow the use of a larger number of packages


% For English and Portuguese languages
% Portuguese will be the default.
% Uncomment \setlanguage below to change it
\usepackage[english,portuguese]{babel}

\usepackage{microtype}  % Make pdf look better

% Side by side figures
% Eg. Fig 1a, Fig 1b
\usepackage[hang,small,bf]{caption}
\usepackage{subcaption}

% Configure Hyperlink color
% As a matter or style, you may use this to enable/disable color boxes on links
%\usepackage[breaklinks=true,colorlinks=false,linkcolor=blue]{hyperref}
% Or use the default values provided by the hyperref package
\usepackage{hyperref}

% Disable commands not supported as PDF string
\pdfstringdefDisableCommands{%
  \def\\{ }%
  \def\texttt#1{#1}%
  \def\selectlanguage#1{}%
}

% Configure pdf information dictionary
\hypersetup{
  pdfauthor={{\ciAuthorFirst} {\ciAuthorLast}},
  pdftitle={\ciTitlePT},
  pdfkeywords={\ciKeywords},
  pdfsubject={\ciAbstractText}
}

% Redefine code boxes
\usepackage{minted}
\makeatletter
\newcommand{\currentfontsize}{\fontsize{\f@size}{\f@baselineskip}\selectfont}
\makeatother

\setminted{fontsize=\footnotesize}
\setmintedinline{fontsize=\currentfontsize}

\addto\captionsportuguese{%
  \renewcommand\listingscaption{Código}
}
\fvset{fontsize=\footnotesize} % Make Code blocks smaller than text
\usepackage{csquotes}

\usepackage{comment} % Add support for PDF Comments
\ifthenelse{\equal{\useoverleaf}{0}}
{\usepackage{pdfcomment}}{}

\usepackage{bookmark} % New Bookmarks
\usepackage{multicol} % For Multiple columns in Glossary
\usepackage{amsmath} % Add support for Math symbols
\usepackage{amssymb} % Add support for Math symbols
\usepackage{eurosym} % Add support for the Euro symbol
\usepackage{todonotes} % Add support for missingfigure and todo

\usepackage{graphicx}
\graphicspath{{figs/}{figures/}}
\usepackage{tikz}
\usetikzlibrary{arrows.meta,positioning,fit,backgrounds,calc,shapes.geometric}
% Shared TikZ styles for architecture block diagrams
\tikzset{
  block/.style  = {draw, rounded corners=2pt, align=center, font=\small,
                    minimum height=2.2em, inner sep=5pt, fill=black!3},
  layer/.style  = {draw, rounded corners=3pt, align=center, font=\small\bfseries,
                    fill=black!6, inner sep=8pt},
  seam/.style   = {draw, dashed, rounded corners=2pt, align=center, font=\small},
  flow/.style   = {-{Stealth[length=2.4mm]}, thick},
  biflow/.style = {{Stealth[length=2.4mm]}-{Stealth[length=2.4mm]}, thick},
}

% Code listings via the pure-LaTeX listings package (no external executable,
% unlike minted v3 which needs latexminted on PATH).
\usepackage{listings}
\lstset{
  basicstyle=\ttfamily\footnotesize,
  breaklines=true,
  columns=fullflexible,
  keepspaces=true,
  showstringspaces=false,
  frame=single,
  framerule=0.3pt,
  rulecolor=\color{black!30},
  xleftmargin=4pt,
  aboveskip=1em,
  belowskip=0.5em,
  captionpos=b,
}
\lstdefinestyle{cpp}{language=C++,
  morekeywords={override,virtual,class,struct,final},
  commentstyle=\itshape\color{black!55}}
\renewcommand{\lstlistingname}{Listing}

% Setup bibliography with Biber using IEEE style for proper UTF-8 support
\usepackage[backend=biber, style=ieee, sorting=none, natbib=true, mincitenames=1, maxcitenames=2]{biblatex}
\bibliography{bib/references.bib}

\usepackage[nohyperlinks]{acronym} % For acronyms (prints every entry in acronyms.tex; the text writes acronyms as plain text, not \ac{})
\usepackage{indentfirst}            % Indenting the first paragraph after section start
\usepackage{xparse}                 % For fixing listoflistings with memoir
\usepackage[detect-weight=true, binary-units=true]{siunitx} % Package to master SI units
\usepackage{verbatim}               % For inline TT text (e.g. code snippets)

%Keep floats inside section!
\usepackage[section]{placeins}
% Barriers only at section/chapter level: a barrier before every subsection
% forced page breaks (and large gaps) after short subsections with big figures.
\let \oldsection \section
\renewcommand{\section}[2][]{
  \FloatBarrier
  \oldsection#1{#2}
}
\let \oldchapter \chapter
\renewcommand{\chapter}[2][]{
  \FloatBarrier
  \oldchapter#1{#2}
}

% extra index for first lines
\makeindex[lines]

\newcommand\mainmatterWithoutReset{
{\edef\temppagenumber{\arabic{page}}%
  \mainmatter
  \setcounter{page}{\temppagenumber + 1}%
}
}

%%%%%%%%%%%%%%%%%%%%%%%%%%%%%%%%%%%%%%%%%%%%%%%%%%
%%       BEGIN DOCUMENT
%%%%%%%%%%%%%%%%%%%%%%%%%%%%%%%%%%%%%%%%%%%%%%%%%%
% Adjust margins according to the Standards for the presentation of the Master Dissertation
%
\settypeblocksize{245mm}{155mm}{*}
\setulmargins{3cm}{*}{*}
\setheadfoot{\onelineskip}{2\onelineskip}
\setlrmarginsandblock{3cm}{2.5cm}{*}
\checkandfixthelayout

% Texts
\newenvironment{introduction}
{%
  \begin{minipage}{\textwidth-2em}%
   \itshape%
}
{%
  \end{minipage}%
  \par\addvspace{2\baselineskip plus 0.2\baselineskip minus 0.2\baselineskip}%
}

% Allow a little extra inter-word stretch so long code identifiers
% (\texttt{...}, which cannot hyphenate) don't run into the margin.
\setlength{\emergencystretch}{3em}

\begin{document}

%%% By default we consider portuguese to be the
%%%  main language (the last one in the usepackage
%%%  declaration). If you're writing your thesis
%%%  in english uncomment the following line.
%%% Set the proper language or you will end with incorrect hyphenation!
\selectlanguage{english}


%% COVER PAGES
% Comment this line to disable cover pages and speed up compilation during editing
\MakeCover
% The cover temporarily uses a full-page (290mm) text block; make sure the
% page-breaking heights are reset to the thesis layout afterwards, otherwise
% long front-matter pages (e.g. the table of contents) run off the page.
\makeatletter
\global\@colht\textheight
\global\@colroom\textheight
\global\vsize\textheight
\makeatother

% Redefine section names according to your preference
%\def\sectionautorefname{Section}
%\def\chapterautorefname{Chapter}
%\def\figureautorefname{Figure}
%\def\listingautorefname{Listing}
%\def\tableautorefname{Table}

% Front matter
% Start front matter, main matter and chapters on the next page (no blank pages for PDF reading)
\let\cleardoublepage\clearpage
\frontmatter

%Exclude sub figures from List of Figures
%\captionsetup[subfloat]{list=no}

% Set chapter style for TOC, LOF, LOT
\pagestyle{plain}
\chapterstyle{veelo}


% Adjust spacing in TOC
\setlength{\cftbeforechapterskip}{0.25em}
\setlength{\cftbeforesectionskip}{0em}
\setlength{\cftbeforesubsectionskip}{0em}
\setlength{\cftbeforesubsubsectionskip}{0em}

\renewcommand\contentsname{Contents}
\small
\SingleSpacing
\tableofcontents
\clearpage
\listoftables
\clearpage
\listoffigures
\clearpage

%\renewcommand{\listingscaption}{Código}
%\renewcommand{\listoflistingscaption}{Lista de Excertos de Código}
%\listoflistings
%\clearpage

% Print Acronyms
\clearpage
\chapter{Acronyms}

\footnotesize
\SingleSpacing

\begin{multicols}{2}
\begin{acronym}[AAAAAA] % The longest acronym length should match the number of A's

	\acro{api}[API]{Application Programming Interface}
	\acro{bim}[BIM]{Building Information Modelling}
	\acro{cad}[CAD]{Computer-Aided Design}
	\acro{cctv}[CCTV]{Closed-Circuit Television}
	\acro{crud}[CRUD]{Create, Read, Update, Delete}
	\acro{dpapi}[DPAPI]{Data Protection Application Programming Interface}
	\acro{dt4mob}[DT4MOB]{Digital Twins for Mobility}
	\acro{ecef}[ECEF]{Earth-Centred, Earth-Fixed}
	\acro{gis}[GIS]{Geographic Information System}
	\acro{gltf}[glTF]{GL Transmission Format}
	\acro{http}[HTTP]{Hypertext Transfer Protocol}
	\acro{hud}[HUD]{Heads-Up Display}
	\acro{ifc}[IFC]{Industry Foundation Classes}
	\acro{iot}[IoT]{Internet of Things}
	\acro{ism}[ISM]{Instanced Static Mesh}
	\acro{json}[JSON]{JavaScript Object Notation}
	\acro{lidar}[LiDAR]{Light Detection and Ranging}
	\acro{oauth}[OAuth]{Open Authorization}
	\acro{ogc}[OGC]{Open Geospatial Consortium}
	\acro{owl}[OWL]{Web Ontology Language}
	\acro{rdf}[RDF]{Resource Description Framework}
	\acro{rest}[REST]{Representational State Transfer}
	\acro{rq}[RQ]{Research Question}
	\acro{rts}[RTS]{Real-Time Strategy (overhead camera mode)}
	\acro{sumo}[SUMO]{Simulation of Urban Mobility}
	\acro{sus}[SUS]{System Usability Scale}
	\acro{ui}[UI]{User Interface}
	\acro{url}[URL]{Uniform Resource Locator}
	\acro{wgs84}[WGS84]{World Geodetic System 1984}

\end{acronym}
\end{multicols}



%%%%%%%%%%%%%%%%%%%%%%%%%%%%%%%%%%%%%%%%%%%%%%%%%%%%%%%
% Main document starts here
%%%%%%%%%%%%%%%%%%%%%%%%%%%%%%%%%%%%%%%%%%%%%%%%%%%%%%%

\mainmatter
\normalsize

\maxtocdepth{subsection} %subsections are added to the TOC
\setsecnumdepth{subsubsection} %subsubsections are numbered

\RenewDocumentCommand \chapter { s O{#3} m }{%
  \FloatBarrier
  \IfValueTF{#1}  % if optional star is seen
    {\oldchapter*{#2}}
    {\oldchapter#1{#2}}
}

% Reset Chapters
\renewcommand{\chapter}[2][]{
  \FloatBarrier
  \oldchapter#1{#2}
}

% Select chapter style from existing or select custom
%\chapterstyle{thesis} % Others: dowding, demo2, dash, chappell, brotherton, bianchi, ger, madsen, tatcher, veelo,indexes)
% Check: https://tug.ctan.org/info/MemoirChapStyles/MemoirChapStyles.pdf

% Default is madsen chapter style


% Set custom section styles
%\setsecheadstyle{\Large\lsstyle\bfseries}
%\setsubsecheadstyle{\large\lsstyle\bfseries}
%\setsubsubsecheadstyle{\normalsize\lsstyle\bfseries}

% Line spacing: 1.5 pt
\OnehalfSpacing

% If you feel adventurous you can also define all aspects of your theme
% Use either this input or the chapterstyle before
% % Rules
\newcommand{\thinRule}{\rule{\textwidth}{0.25pt}}

% Customize heading appearances
% Define styles
\newcommand{\partSize}{\Huge}
\newcommand{\partStyle}{\lsstyle\scshape}
\newcommand{\chapterSize}{\Huge}
\newcommand{\chapterStyle}{\lsstyle\scshape}
\newcommand{\chapterAfter}{}
\newcommand{\sectionSize}{\Large}
\newcommand{\sectionStyle}{\scshape\MakeTextLowercase}
\newcommand{\subsectionSize}{\large}
\newcommand{\subsectionStyle}{\scshape\MakeTextLowercase}
\newcommand{\subsubsectionSize}{\large}
\newcommand{\subsubsectionStyle}{\scshape\MakeTextLowercase}
\newlength{\partNumSizePt}
\setlength{\partNumSizePt}{60pt}
\newlength{\chapterNumSizePt}
\setlength{\chapterNumSizePt}{60pt}
\newcommand{\partNumSize}{%
  \fontsize{\partNumSizePt}{1.2\partNumSizePt}\selectfont%
}
\newcommand{\partNumStyle}{\partChapterNumColor}
\newcommand{\chapterNumSize}{%
  \fontsize{\chapterNumSizePt}{1.2\chapterNumSizePt}\selectfont%
}
\newcommand{\chapterNumStyle}{\partChapterNumColor}

% Customize parts
\renewcommand{\partnamefont}{\partSize\partStyle}
\renewcommand{\partnumfont}{\partNumSize\partNumStyle}
\renewcommand{\printpartname}{}
\renewcommand{\printparttitle}[1]{%
  \normalfont\normalcolor\partnamefont #1
}

% Customize chapters
\makeatletter
\setlength{\beforechapskip}{30pt}
\renewcommand*{\chapterheadstart}{\vspace*{\beforechapskip}}
\setlength{\afterchapskip}{3ex}
\setlength{\midchapskip}{3ex}
\renewcommand*{\chapnamefont}{%
  \Large\flushright\chapterStyle\partChapterNumColor%
}
\renewcommand*{\chapnumfont}{\chapterNumSize\chapterNumStyle}
\renewcommand*{\chaptitlefont}{%
  \normalfont\flushleft\normalcolor\chapterSize\chapterStyle%
}
\renewcommand*{\printchaptername}{%
  \chapnamefont\MakeTextLowercase{\@chapapp}%
}
\renewcommand*{\chapternamenum}{\quad}
\renewcommand*{\printchapternum}{%
%  \chapnumfont\textls[-75]{\classicstylenums{\thechapter}}%
 \chapnumfont\textls[-75]{\thechapter}%

}
\renewcommand*{\printchaptertitle}[1]{%
  \chaptitlefont #1
  \chapterAfter
}
\makeatother
% Customize sections and subsections
\setsecnumformat{\csname my#1\endcsname\quad}
\setsecheadstyle{\sectionSize\sectionStyle}
\newcommand{\mysection}{{\thesection}}
\setlength{\beforesecskip}{3em}


\setsubsecheadstyle{\subsectionSize\subsectionStyle}
\newcommand{\mysubsection}{{\normalfont\subsectionSize\thesubsection}}
\setlength{\beforesubsecskip}{3em}

\setsubsubsecheadstyle{\subsubsectionSize\subsubsectionStyle}
\newcommand{\mysubsubsection}{{\normalfont\subsubsectionSize\thesubsubsection}}
\setlength{\beforesubsubsecskip}{2em}

% Customize "Table of ..." appearance
% Customize headings
\newcommand{\renewPrintXTitle}[1]{%
  \renewcommand{#1}[1]{%
    \printchaptertitle{##1}%
  }%
}
\renewPrintXTitle{\printtoctitle}
\renewPrintXTitle{\printlottitle}
\renewPrintXTitle{\printloftitle}

% Customize ToC headings
\renewcommand{\cftpartfont}{\partChapterNumColor\partStyle}
\renewcommand{\cftchapterfont}{\chapterStyle}
\renewcommand{\cftsectionfont}{}
\renewcommand{\cftsubsectionfont}{}
\renewcommand{\cftfigurefont}{}
\renewcommand{\cfttablefont}{}
\newcommand{\cftlstlistingfont}{}

% Increase number width
\newlength{\cftNumWidthIncrease}
\setlength{\cftNumWidthIncrease}{0.25em}
\addtolength{\cftpartnumwidth}{\cftNumWidthIncrease}
\addtolength{\cftchapternumwidth}{\cftNumWidthIncrease}
\addtolength{\cftsectionindent}{\cftNumWidthIncrease}
\addtolength{\cftsubsectionindent}{\cftNumWidthIncrease}
% No leader dots
%\renewcommand*{\cftpartdotsep}{\cftnodots}
%\renewcommand*{\cftchapterdotsep}{\cftnodots}
%\renewcommand*{\cftsectiondotsep}{\cftnodots}
%\renewcommand*{\cftsubsectiondotsep}{\cftnodots}
%\renewcommand*{\cftfiguredotsep}{\cftnodots}
%\renewcommand*{\cfttabledotsep}{\cftnodots}
%\newcommand*{\cftlstlistingdotsep}{\cftnodots}
% Set page numbers immediately after entry text
\newcommand{\tocEntryPageSep}{\hspace{1em}}
\renewcommand{\cftpartleader}{\cftdotfill{\cftdotsep}}
%\renewcommand{\cftpartafterpnum}{\cftparfillskip}
%\renewcommand{\cftchapterleader}{\tocEntryPageSep}
\renewcommand{\cftchapterleader}{\cftdotfill{\cftdotsep}}
%\renewcommand{\cftchapterafterpnum}{\cftparfillskip}
\renewcommand{\cftsectionleader}{\cftdotfill{\cftdotsep}}
%\renewcommand{\cftsectionafterpnum}{\cftparfillskip}
\renewcommand{\cftsubsectionleader}{\cftdotfill{\cftdotsep}}
%\renewcommand{\cftsubsectionafterpnum}{\cftparfillskip}
\renewcommand{\cftfigureleader}{\cftdotfill{\cftdotsep}}
%\renewcommand{\cftfigureafterpnum}{\cftparfillskip}
\renewcommand{\cfttableleader}{\cftdotfill{\cftdotsep}}
%\renewcommand{\cfttableafterpnum}{\cftparfillskip}
\newcommand{\cftlstlistingleader}{\cftdotfill{\cftdotsep}}
%\newcommand{\cftlstlistingafterpnum}{\cftparfillskip}
% Customize page numbers
\newcommand{\tocPageStyle}{\tocPageColor}
\renewcommand{\cftpartpagefont}{\tocPageStyle}
\renewcommand{\cftchapterpagefont}{\tocPageStyle}
\renewcommand{\cftsectionpagefont}{\tocPageStyle}
\renewcommand{\cftsubsectionpagefont}{\tocPageStyle}
\renewcommand{\cftfigurepagefont}{\tocPageStyle}
\renewcommand{\cfttablepagefont}{\tocPageStyle}
\newcommand{\cftlstlistingpagefont}{\tocPageStyle}

% Abstract
% Remove indents around abstract text
\setlength{\absleftindent}{0pt}
\setlength{\absrightindent}{0pt}
% Change font size to conform with the rest of the document text
\renewcommand{\abstracttextfont}{\normalsize}

% Customize headers and footers including page numbers
\newcommand{\hfTextSize}{\footnotesize}
\newcommand{\headTextStyle}{\lsstyle\scshape\MakeTextLowercase}
\nouppercaseheads
\makeevenhead{headings}%
             {\hfTextSize\thepage}%
             {}%
             {\hfTextSize\headTextStyle\leftmark}
\makeevenhead{plain}%
             {\hfTextSize\thepage}%
             {}%
             {\hfTextSize\headTextStyle\leftmark}
\makeoddhead{headings}%
            {\hfTextSize\headTextStyle\rightmark}%
            {}%
            {\hfTextSize\thepage}
\makeoddhead{plain}%
            {\hfTextSize\headTextStyle\rightmark}%
            {}%
            {\hfTextSize\thepage}


% Customize captions
\newcommand{\captionSize}{\small}
\newcommand{\captionStyle}{\scshape}
\newcommand{\captionWidthRatio}{0.9}

\captionnamefont{\captionSize\captionStyle}
\captiontitlefont{\captionSize}
\captiondelim{ -- }
\captiontitlefinal{}
\changecaptionwidth
%\captionwidth{\captionWidthRatio\textwidth}

% Define colors
%\newcommand{\titleColor}{\color[rgb]{0.616, 0.0627, 0.176}}
\newcommand{\titleColor}{\color[rgb]{0,0,0}}

\newcommand{\partChapterNumColor}{\titleColor}
\newcommand{\dropCapColor}{\titleColor}
%\newcommand{\tocPageColor}{\color[rgb]{0.0980, 0.329, 0.651}}

\newcommand{\tocPageColor}{\color[rgb]{0, 0,0}}
\definecolor{shade0}{rgb}{1.0 , 1.0 , 1.0 }
\definecolor{shade1}{rgb}{0.9 , 0.9 , 0.9 }
\definecolor{shade2}{rgb}{0.8 , 0.8 , 0.8 }
\definecolor{shade3}{rgb}{0.65, 0.65, 0.65}
\definecolor{shade4}{rgb}{0.45, 0.45, 0.45}
\definecolor{shade5}{rgb}{0.0 , 0.0 , 0.0 }



%% CHAPTERS
\chapter{Introduction}
\label{chap:introduction}

% TODO: draft, expand once the rest of the thesis exists
\begin{introduction}
This chapter introduces the problem addressed by this thesis, states its objectives and contributions, and outlines the structure of the remaining chapters.
\end{introduction}

\section{Context and Motivation}
% TODO: digital twin systems for physical infrastructure are typically built as
% closed, purpose-specific applications, tightly coupled to one infrastructure
% type and one data backend. Motivate the need for a generalized, decoupled
% viewer.

\section{Thesis Objectives}
% TODO: state the objectives explicitly, e.g.:
% - design and implement an open-source digital twin viewer generalizable
%   across infrastructure types
% - decouple the viewer's rendering/interaction layer from any single data
%   backend
% - evaluate the resulting architecture through performance benchmarks and
%   user testing

\section{Contributions}
% TODO: list concrete contributions (the viewer itself as open-source
% software, the architectural pattern, the benchmark harness/results, the
% user study once completed)
\begin{itemize}
    \item % TODO
\end{itemize}

\section{Thesis Structure}
This document is organized into six chapters. Chapter~\ref{chap:introduction} introduces the problem and objectives. Chapter~\ref{chap:state-of-the-art} reviews related work on digital twin visualization and geospatial engines. Chapter~\ref{chap:architecture} presents the high-level architecture of the viewer. Chapter~\ref{chap:implementation} details the implementation decisions made to achieve generalization and decoupling. Chapter~\ref{chap:evaluation} evaluates the resulting system through performance benchmarks and user testing. Chapter~\ref{chap:conclusions} summarizes the contributions and discusses limitations and future work.

\chapter{State of the Art}
\label{chap:state-of-the-art}

\begin{introduction}
This chapter reviews the background and existing work relevant to the viewer developed in this thesis. It states how the literature was gathered, then covers the application of digital twins, introduced in Chapter~\ref{chap:introduction}, to physical infrastructure, the data sources and models that infrastructure digital twins draw on, the shift from dedicated Geographic Information System (GIS) tooling to real-time game engines for geospatial visualisation, and the handling of simulated entities along with physical ones. It then examines existing digital twin platforms and viewers, and how tightly each couples its rendering to a particular backend and data model. The chapter closes by stating the gap this thesis targets.
\end{introduction}

\section{Review Methodology}
\label{sec:review-methodology}

The literature underlying this chapter, and the digital twin background in Chapter~\ref{chap:introduction}, was gathered through a keyword search across Scopus, IEEE Xplore, Google Scholar, and ScienceDirect, using broad term combinations such as ``digital twin'', ``digital twin'' AND ``game engine'', and ``digital twins in game engines''. Because ``digital twin'' spans disciplines from manufacturing to healthcare to civil engineering, these broad terms returned a large, heterogeneous candidate pool; the Scopus search alone returned 320 records. This pool was narrowed by title and abstract relevance to the intersection this thesis actually addresses, meaning infrastructure digital twins, real-time game-engine visualisation, and generalisable viewer architectures, discarding the substantial majority of results that fell in unrelated domains (manufacturing digital twins, medical imaging, structural digital twins with no visualisation component, and similar).

The directly relevant set that remained after this narrowing is small, reflecting how narrow this specific intersection is in the published literature rather than an artefact of the search. Some of the works ultimately cited, including the foundational digital twin papers by \textcite{grieves2016digitaltwin} and \textcite{glaessgen2012digitaltwin}, were identified directly as known origin points for the concept rather than surfaced by the keyword search; further sources, including \textcite{broo2021digitaltwins}, \textcite{kherbache2022ditto}, and the geospatial-visualisation and related-viewer-system literature in Sections~\ref{sec:geospatial-vis} and~\ref{sec:related-viewer-systems}, were found by snowballing: following the reference lists and citing works of the papers already found this way.

\section{Digital Twins for Physical Infrastructure}
\label{sec:dt-infrastructure}

\textcite{broo2021digitaltwins} narrow the digital twin concept, introduced in Section~\ref{sec:context}, to infrastructure specifically, describing digital twins in this domain as an emerging idea without a commonly accepted definition, but converging on a realistic digital representation of assets, processes, or systems in the built environment, kept synchronised with their physical counterpart. In the mobility sub-domain, entities are the components of a road network and the sensing and telematics equipment that instruments them, as described in Section~\ref{sec:context}, and digital twins are expected to support applications such as traffic management, predictive maintenance, incident detection, and scenario evaluation~\cite{fuller2020digitaltwin}.

Architecturally, an infrastructure digital twin is conventionally split into a backend that models and serves entity data and a viewer that renders it. The backend's role is to abstract over the heterogeneity of the physical layer, with its many device types and protocols, and expose a uniform interface. Eclipse Ditto is open-source middleware built for exactly this: it represents physical devices as digital twins behind a common HTTP and WebSocket API, decoupling client applications from how each device actually communicates~\cite{kherbache2022ditto}. Ditto is explicitly a data layer; it provides no visualisation. \textcite{kherbache2022ditto} pair it with Eclipse Hono for device connectivity and demonstrate the combination as a digital twin network for the industrial IoT. Still, any three-dimensional or geospatial representation of the twins it manages must be built separately.

This backend/viewer separation is near-universal in the surveyed literature: the system that models the twin's data and the system that renders it are treated, and built, as distinct concerns. Where systems fall short of that separation \emph{within} the viewer, meaning the viewer's own code is written against one backend and one entity schema, is the subject of Section~\ref{sec:related-viewer-systems}.

This coupling leaves a gap that is not simply anecdotal. \textcite{masoumi2023citydigitaltwins} systematically reviewed 41 published city digital twin studies and found that 90\% of them sit at an initial-to-medium maturity level, concentrated on 3D modelling, monitoring, and visualisation as opposed to genuine two-way interaction; they identify near-real-time data analytics and improved BIM/GIS integration as open needs rather than solved problems. This independently corroborates, across a much larger sample than this chapter surveys directly, that the field as a whole has not moved much past the coupled default described above.

\section{Data Sources and Models for Infrastructure Digital Twins}
\label{sec:dt-data}

An infrastructure digital twin draws on several kinds of data, each with its own structure and rate of change. Producing and storing that data is the backend's concern rather than the viewer's (Section~\ref{sec:scope}), but the kinds of data involved determine what the viewer has to display: semantically rich models of static assets, continuous streams of sensor readings, and large numbers of moving vehicles.

\subsection{Building Information Modelling}

Building Information Modelling (BIM) is a methodology for creating and managing digital representations of the physical and functional characteristics of built assets, with detailed geometric, semantic, and relational information about each element, such as its dimensions, materials, and lifecycle attributes~\cite{azhar2011bim}. Although BIM originated in the architecture, engineering, and construction sector, it is increasingly applied to large infrastructure such as roads, bridges, and tunnels. Within digital twin frameworks, BIM models provide a high-fidelity structural and semantic layer; enriched with operational and real-time data, a BIM model stops being a static design artefact and becomes an active part of a digital twin, supporting asset management, condition monitoring, and maintenance planning~\cite{deng2021bimdigitaltwins}. For a viewer, such models are the source of asset footprints and attributes, such as those in a road operator's inventory of signage, barriers, and pavement sections.

\subsection{Sensors and Data Acquisition}

Sensors measure properties of the environment or of a structure, such as traffic flow, vehicle speed, weather conditions, structural stress, or air quality. In intelligent transportation systems they are the primary interface between the physical infrastructure and digital platforms, providing continuous data from the road network~\cite{felemban2014its}. Common types include inductive loop detectors, cameras, radar, LiDAR, environmental sensors, and structural health monitoring devices. The data they produce is heterogeneous, frequent, and geographically distributed, which makes it hard to integrate, process, and visualise. For a viewer, each sensor is an entity whose state changes live, such as an inclinometer on a retaining wall or a roadside traffic counter.

\subsection{Connected Vehicles}

Connected vehicles exchange data with other vehicles (vehicle-to-vehicle, V2V), with roadside infrastructure (vehicle-to-infrastructure, V2I), and with networked systems (vehicle-to-network, V2N), collectively referred to as vehicle-to-everything (V2X) communication~\cite{abboud2016v2x}. They report their position, speed, and trajectory, as well as observed road conditions and traffic events, and integrating this information into a digital twin of a road corridor gives a detailed, up-to-date picture of traffic~\cite{wu2025v2xdigitaltwin}. For a viewer, connected vehicles mean large numbers of moving entities whose positions change with every update, similar to the simulated vehicles discussed in Section~\ref{sec:simulated-entities}.

\section{Geospatial Visualisation: From GIS to Game Engines}
\label{sec:geospatial-vis}

Rendering a digital twin of physical infrastructure at real-world scale requires placing content accurately on a geographic reference frame. This has historically been the province of GIS software~\cite{goodchild2010giscience} and of standardised 3D city models: CityGML and its JSON encoding, CityJSON, describe urban objects, such as buildings, roads, vegetation, and city furniture, with defined geometry and semantics, and most urban digital twin platforms are built around such a model. The bridge from a BIM-side model to this representation is itself an active area: \textcite{xu2019ifc3dtiles} convert Industry Foundation Classes (IFC) building models into the 3D Tiles format this thesis also streams through, and \textcite{lam2024digitaltwinsmartcity} similarly map IFC into CityGML and an OWL/RDF semantic graph, visualising the result through both Cesium Ion and Unreal Engine; both treat 3D Tiles and CityGML as a common target format for otherwise incompatible authoring tools rather than something specific to one project.

Real-time 3D game engines have become a practical alternative as they have gained geospatial capabilities. \textcite{rantanen2023opengeospatial} integrate open geospatial data into a game engine for urban digital twin applications and describe the data-conversion pipeline this requires; \textcite{ene2026urbandigitaltwins} survey the combined use of GIS and game engine systems for urban digital twins and find game engines increasingly viable where interactivity and rendering quality matter more than analytical GIS functionality. Both note the trade-off: game engines bring strong rendering and interaction but must be bridged to geospatial data and coordinate systems they do not natively understand. \textcite{virtanen2020pointclouds} push this further for raw sensor data, extending a game engine to interactively render dense coloured point clouds instead of only pre-processed city models.

The coordinate systems involved recur throughout this thesis and are worth stating. Geographic positions are almost always expressed in the World Geodetic System 1984 (WGS84), the global reference system used by GPS, which models the Earth as an ellipsoid and identifies each point by its latitude, longitude, and height above that ellipsoid~\cite{nima2000wgs84}. Latitude and longitude are angles rather than distances, so they cannot be used directly as positions in a game engine, which works in a flat Cartesian space measured in fixed units. A common intermediate step is the Earth-Centred, Earth-Fixed (ECEF) frame, a Cartesian system in metres whose origin is the Earth's centre of mass and which rotates with the Earth. A WGS84 position converts to ECEF with a fixed formula, and ECEF positions can then be expressed relative to a local origin near the area being viewed, which yields the flat, metre-scale coordinates an engine expects.

The bridge used in this work is Cesium for Unreal~\cite{cesiumforunreal}, a free and open-source plugin that embeds a full-scale WGS84 globe in Unreal Engine and streams terrain, imagery, and 3D content into it. The streaming format is 3D Tiles, an OGC Community Standard for organising and delivering massive heterogeneous 3D geospatial datasets, such as point clouds, photogrammetry, BIM and CAD models, and 3D buildings, as a spatially indexed hierarchy suitable for progressive loading~\cite{ogc3dtiles}. Cesium for Unreal was chosen over building geospatial placement from scratch, and over a pure web or GIS viewer, because it retains Unreal Engine's rendering and interaction capabilities, relevant for a viewer intended to generalise across infrastructure types with potentially rich per-entity visuals, while still anchoring content to accurate real-world coordinates. Section~\ref{sec:geospatial-placement} in Chapter~\ref{chap:architecture} describes how this thesis builds a single coordinate-conversion seam on top of Cesium's georeferencing rather than scattering coordinate mathematics across the application.

\section{Simulated Entities Along with Physical Ones}
\label{sec:simulated-entities}

A viewer for a mobility digital twin is increasingly expected to display not only physical infrastructure but also entities produced by simulation and streamed through the same backend. Two cases arose in this work. The first is traffic: microscopic traffic simulators such as SUMO model individual vehicles' movement through a road network~\cite{lopez2018sumo}, producing a continuous, high-churn stream of short-lived vehicle entities, each with a position that changes with every update. The second is hazard modelling: a wildfire-spread simulation produces a projected burn perimeter and spread cone that evolve and are attached to an ignition-point entity.

These entities differ from physical infrastructure in ways that matter to the viewer. They are numerous and ephemeral, update far more frequently, may share a single fleet-wide lifetime rather than expiring individually, and may have their geometry delivered as GeoJSON rather than as a mesh. A viewer whose entity handling is specialised to static infrastructure will not accommodate them without change. This reinforces the need, argued throughout this thesis, for an entity model that is generic over what an entity is rather than assuming one shape.

A related case from outside the mobility domain is \textcite{yang2023cesiumlandslide}'s web-based emergency-response system, also built on Cesium Digital Earth, which maps live IoT and sensor-network data for rescue workers, vehicles, and monitored landslide sites onto the globe for search-and-rescue coordination. It is a hazard-adjacent parallel to the wildfire case above, streaming heterogeneous physical entities (personnel, equipment, instrumented sites) through one viewer. However, the paper does not report whether its entity handling generalises beyond this fixed set, as this thesis argues it should.

\section{Digital Twin Platforms and Related Viewer Systems}
\label{sec:related-viewer-systems}

Digital twin platforms for cities and mobility are now numerous, and a recent systematic review maps how IoT, artificial intelligence, and digital twins are combined across smart-city systems~\cite{sacoto2025smartcities}. The systems discussed below are those closest to the viewer in this thesis: each pairs a source of data with a 3D or geospatial view of infrastructure or mobility, and together they are useful comparators for how tightly the view is coupled to the data behind it.

City2Twin~\cite{rafamatanantsoa2024city2twin} presents an open urban digital twin pipeline that integrates municipal open data into a 3D city model for visualisation and analysis. Like most urban digital twin platforms surveyed by \textcite{ene2026urbandigitaltwins}, it is organised around a specific data pipeline and city-model format, typically CityGML or CityJSON. It is not designed to be repointed at an arbitrary IoT backend or entity schema. The 3D city model is both its input and its internal representation.

Cesium for Unreal itself, and the wider Cesium platform, are backend-agnostic in principle: they render whatever geospatial data they are given. But they ship no entity model, no live-update mechanism, and no interaction layer for IoT-style digital twins. Applications built on Cesium, including those surveyed in Section~\ref{sec:geospatial-vis}, still implement their own data layer, and that layer is typically written for one project's backend and one project's entity types. The geospatial foundation is reusable; the digital twin logic on top of it generally is not.

The closest domain match in the surveyed literature is \textcite{lv2024highwaysandbox}'s digital sandbox platform, also built on Unreal Engine, for highway route planning and design. It integrates BIM data from a fixed commercial toolchain (Bentley OpenRoads Designer and OpenBridge Modeler) with a photogrammetry-derived GIS terrain model into what the authors call a digital base model, then builds route-selection, traffic-simulation, and component-query modules on top of it. Two differences from the viewer in this thesis are instructive: it is a design-time planning tool built from a static, pre-processed model rather than a viewer kept current by a live backend connection, and its component types, meaning bridges, tunnels, and interchanges, are fixed to that one BIM/GIS pipeline rather than expressed through a generic, extensible entity model. Even a viewer this close in domain and engine choice does not attempt either of the two axes of generalisation this thesis argues for.

CARLA~\cite{dosovitskiy2017carla}, an open-source urban driving simulator built on Unreal Engine, shows how far a game engine can be taken in the mobility domain: it provides high-fidelity 3D urban environments, realistic traffic scenarios, and simulated sensors such as cameras and LiDAR. Its purpose, however, is developing and testing autonomous-driving systems at the level of individual vehicles, not monitoring a real road network through a live connection to infrastructure data. Closer to the connected road corridors that mobility digital twins target, \textcite{wu2025v2xdigitaltwin} propose a digital twin framework for integrating the physical and virtual sides of V2X-enabled connected-vehicle corridors; its contribution is that integration framework, and it does not address a viewer that generalises across entity types or backends.

Outside the geospatial domain, \textcite{robles2023opentwins}'s OpenTwins is an open-source framework for \emph{compositional} digital twins, in which individual entities and subsystems are linked and orchestrated with 3D visualisation, IoT data streams, and machine-learning predictions. Its composition-over-monolith stance is close in spirit to the entity model in Chapter~\ref{chap:architecture}, and it demonstrates that open, generalisable digital twin tooling is achievable; its target domain, however, is Industry 4.0 process orchestration, not geospatial mobility infrastructure, and it does not address placing entities on a real-world globe or streaming them by geographic area.

The academic and open-source systems above have commercial, closed-source counterparts, and these are worth naming because they show the coupled default is not a limitation of small or under-resourced projects: it persists in mature, widely deployed enterprise tooling. Bentley's iTwin platform~\cite{bentleyitwin} is a cloud digital twin ecosystem built around \emph{iModels}, Bentley's own proprietary schema (the Base Infrastructure Schema, BIS), populated primarily from Bentley's own authoring tools, such as MicroStation, OpenRoads, and OpenBuildings, the same product family used for the BIM side of the highway sandbox platform above. Siemens' Xcelerator portfolio~\cite{siemensxcelerator} plays a similar role for industrial and infrastructure IoT, connecting a digital twin to sensor data through Siemens' own device and data-model ecosystem. Esri's ArcGIS platform~\cite{esriarcgis}, extended for city- and infrastructure-scale digital twins through products such as ArcGIS Urban, takes the same shape from the GIS side: entities and their attributes are modelled as Esri geodatabase feature classes, and the surrounding tooling, analysis, and visualisation are all built to that one format. In each case, the viewer, the data model, and often the authoring tools are a single vendor-integrated product; none is designed to point at an arbitrary third-party backend or accept an entity type it was not built to know about, which is precisely the coupling this thesis's architecture avoids.

Across these systems a common pattern appears: the viewer's rendering and interaction code is written against one backend's data model, one fixed set of entity types, or both. This is a defensible choice for a single-deployment viewer, but it means the viewer cannot be reused for a different infrastructure domain or backend without editing its core, because no seam separates ``what an entity looks like and does'' or ``where its data comes from'' from the rest of the application.

\section{Summary and Research Gaps}
\label{sec:research-gaps}

The literature establishes the digital twin as a physical entity, a virtual representation, and a live connection between them. It establishes infrastructure digital twins as a backend/viewer split in which the backend (for example, Eclipse Ditto) models entity data and the viewer renders it. It shows that game engines paired with geospatial plugins such as Cesium for Unreal, streaming standardised formats such as 3D Tiles, are now a viable alternative to GIS-specific tooling for building that viewer. It also shows that such viewers are increasingly expected to handle very different kinds of data at once: semantically rich asset models, live sensor readings, moving vehicles, and simulated entities such as traffic and hazards.

What the surveyed viewer systems do not generally address is generalisation and decoupling \emph{within} the viewer. Existing systems tend to fix both the entity types they can display and the backend they fetch from, so neither can change without changing the viewer's core code, and the viewer cannot run without a live instance of its specific backend. The gap matters most where one viewer must serve several infrastructure domains at once, since a viewer fixed to one of them cannot serve the others. This thesis targets that gap: it presents a viewer whose entity model and data-access layer are each a single replaceable seam rather than assumptions baked throughout the rendering and interaction code, as detailed in Chapter~\ref{chap:architecture}.

\chapter{Architecture}
\label{chap:architecture}

% TODO: draft, expand with diagrams
\begin{introduction}
This chapter presents the high-level architecture of the digital twin viewer, focusing on the design choices that make it generalizable across infrastructure types and decoupled from any single data backend.
\end{introduction}

\section{Overview}
% TODO: high-level component diagram — service layer, entity model,
% coordinate conversion, camera/interaction layer, UI

\section{Service Layer}
% TODO: GameInstance subsystem services as the boundary between the app and
% external data sources; why this boundary enables backend-agnosticism

\section{Entity Model and Generalization}
% TODO: entity struct schemas + factory pattern (thing-type -> struct
% mapping) as the mechanism for supporting arbitrary infrastructure types
% without touching core rendering/interaction code

\section{Geospatial Placement}
% TODO: coordinate conversion service, why all placement goes through one
% seam rather than being scattered across actors

\section{Design Rationale}
% TODO: articulate the architectural principle underlying the above
% (decoupling data model from rendering, rendering from backend) and how it
% differs from the closed/coupled systems surveyed in Chapter~\ref{chap:state-of-the-art}

\chapter{Implementation}
\label{chap:implementation}

\begin{introduction}
This chapter walks through the concrete implementation decisions made throughout development to realize the architecture presented in Chapter~\ref{chap:architecture}, using them as case studies of the architecture's generalization and decoupling in practice. Each section takes one problem encountered during development, describes how it was solved, and relates the solution back to the seam it belongs to. Where a decision was later revised, that revision is described as well, since the revisions are themselves evidence of where the boundaries did and did not hold.
\end{introduction}

\section{Overview}
\label{sec:impl-overview}

The viewer did not start out decoupled. Its earliest form assumed a single backend, read its credentials from a checked-in configuration file, spawned a fixed set of actor types with type-specific code in the spawn path, and carried fire-simulation-specific logic in the generic entity code. The structure described in Chapter~\ref{chap:architecture} -- the service layer, the registration table, the two extension points, the coordinate seam -- is the result of a sequence of refactorings, each triggered by a concrete need: a second way of supplying credentials, a new entity type that did not fit the existing pattern, a rendering feature that only one entity type needed.

The sections below follow that sequence. Section~\ref{sec:impl-backend} covers moving backend specifics out of code. Section~\ref{sec:impl-streaming} covers streaming entities by geographic area and the reliability work that entailed. Section~\ref{sec:impl-worked-example} works through adding a new entity type end to end. Section~\ref{sec:impl-camera} covers the camera and interaction layer and why it was unaffected by the entity-side changes. Section~\ref{sec:impl-special} covers the one place genuinely domain-specific behaviour lives, and how it was kept behind the same seam as every other type. Section~\ref{sec:impl-discussion} reflects on which decisions generalised well and which introduced complexity.

\section{Decoupling the Data Backend}
\label{sec:impl-backend}

Three related changes moved knowledge of the backend out of the application's code and into configuration and runtime state.

\subsection{From a Configuration File to a Data Asset}
\label{sec:impl-dataasset}

Credentials and connection settings were originally read from a \texttt{Secrets.ini} file, git-ignored, parsed directly wherever they were needed. This had two problems: the file was parsed independently by more than one service, so their interpretations could drift, and configuration was expressed as text to be parsed rather than as structured data.

The first step replaced the file with a \texttt{UDittoSecretsAsset} data asset -- a first-class engine object, edited in the editor, loaded exactly once in \texttt{UDittoService::Initialize} from a fixed content path. This made connection settings structured data with a single load point. It is worth being precise about what this did and did not achieve: an early commit described the asset as ``encrypted'', but cooking a data asset into a package does not by itself encrypt it; real protection requires enabling package encryption and AES keys in the project's cryptography settings. The honest description is ``configuration as data, loaded once'', not ``encrypted secrets''.

\subsection{Runtime Credentials and the Startup Sequence}
\label{sec:impl-auth}

The second step moved credentials out of design-time configuration entirely. In the current design, the data asset holds only non-secret defaults (host, protocol, OAuth client id, the stream-start message) for developer convenience. Real credentials arrive at runtime: the user enters them on a dedicated login screen, and \texttt{UDittoService::Login} performs an OAuth2 password grant against a Keycloak endpoint. \texttt{Initialize} no longer authenticates at all.

The consequence is that authentication is asynchronous relative to the rest of startup, and several components need the result. This is coordinated entirely through delegates and a request queue, shown in Figure~\ref{fig:auth-sequence}: any REST call issued before a token lands is queued inside \texttt{UDittoService} and flushed automatically once it does, and \texttt{UWSService} does not open its socket until \texttt{UDittoService} broadcasts that an authorisation header is available. No caller has to reason about login timing; it either gets a live result or a result that arrives slightly later.

\begin{figure}[t]
  \centering
  \begin{tikzpicture}[node distance=6mm, every node/.style={font=\small}]
    \node[block, text width=26mm] (login) {login screen};
    \node[block, right=14mm of login, text width=30mm] (ditto) {\texttt{UDittoService}};
    \node[block, right=14mm of ditto, text width=24mm] (kc) {Keycloak};
    \node[block, below=12mm of ditto, text width=30mm] (ws) {\texttt{UWSService}};
    \node[block, below=8mm of ws, text width=30mm] (daemon) {\texttt{UEntityUpdateDaemon}};

    \draw[flow] (login) -- node[above,font=\footnotesize]{Login()} (ditto);
    \draw[flow] (ditto) -- node[above,font=\footnotesize]{password grant} (kc);
    \draw[flow] (kc) -- node[below,font=\footnotesize]{access + refresh token} (ditto);
    \draw[flow] (ditto) -- node[right,font=\footnotesize]{OnAuthHeaderReady} (ws);
    \draw[flow] (ws) -- node[right,font=\footnotesize]{connect + START-SEND-EVENTS} (daemon);
    \node[right=1mm of ditto.north east, font=\footnotesize\itshape, anchor=south east, yshift=6mm] {queued REST calls flushed on token};
  \end{tikzpicture}
  \caption[The authentication and startup sequence]{Authentication is asynchronous relative to startup. Components subscribe to \texttt{UDittoService}'s ``header ready'' broadcast; REST calls made before the token arrives are queued and replayed, so no caller handles login timing explicitly.}
  \label{fig:auth-sequence}
\end{figure}

Persistence of credentials is handled by a separate subsystem, \texttt{UCredentialStoreService}, backed by platform secure storage (Windows DPAPI or macOS Keychain Services). Only the host, username, and password are persisted; the short-lived OAuth access and refresh tokens are deliberately never written to disk. Instead, a fresh password grant is performed on every launch. This is a deliberate trade: it avoids the failure mode where a persisted refresh token silently goes stale between sessions, at the cost of one extra round-trip at startup, and it means a stored login stays usable until the user explicitly logs out. On a platform with no secure-storage backend implemented, loading simply fails and the login screen is shown every launch -- a graceful degradation rather than an error.

The login flow is isolated at the level design as well: it runs under its own minimal game mode with no pawn, no world geometry, and no Cesium globe, so none of the tile-streaming machinery starts before the user is authenticated.

\subsection{A Single Source of Truth for Connection Settings}

With settings loaded once by \texttt{UDittoService}, the other subsystems that need them read them from its getters rather than loading the asset a second time. \texttt{UWSService} in particular takes its host, protocol, and stream-start message this way. The result is that there is one interpretation of the connection configuration in the running application, and one moment at which the live connection is considered ready, rather than each subsystem arriving at its own.

\subsection{What Backend Replacement Would Still Require}
\label{sec:impl-replace-caveat}

It is important to be precise about the sense in which the backend is ``replaceable''. What has been achieved is \emph{confinement}: no code above the service layer names an HTTP verb, a URL, an OAuth scope, a WebSocket frame, or Ditto's envelope format, so the blast radius of a backend change is one layer. What has \emph{not} been done is to factor \texttt{UDittoService} behind an abstract interface, or to build a second backend implementation or an offline mock. A genuine swap today would mean rewriting the internals of \texttt{UDittoService}, \texttt{UWSService}, and the envelope parser in \texttt{UEntityUpdateDaemon} against the new backend, while keeping their public signatures. That is a bounded, well-localised task -- which is the point of the seam -- but it is not yet a drop-in replacement. This is revisited in Chapter~\ref{chap:conclusions}.

\section{Entity Streaming and Reliability}
\label{sec:impl-streaming}

\subsection{Streaming Entities by Geographic Area}

Loading every entity in the dataset regardless of where the camera is looking does not scale, and baking in an assumption about dataset size would undermine the goal of generalising across infrastructure of arbitrary extent. Instead, \texttt{ADT4MOBGamemode} maps the camera's altitude to a discrete zoom level and, above a threshold zoom, fetches only the entities within the 3$\times$3 grid of geographic tiles around the camera, expressed as quadkey ranges and queried through \texttt{UDittoService::GetThingsByGeotileBounds}. Below the threshold zoom, tile filtering is disabled and everything loads, since at that range the whole area of interest is on screen anyway. Figure~\ref{fig:tile-flow} shows the per-check decision path.

\begin{figure}[t]
  \centering
  \begin{tikzpicture}[node distance=5mm, every node/.style={font=\small}]
    \node[block, text width=40mm] (raw) {camera $\rightarrow$ lat/lon/alt\\{\footnotesize raw zoom = f(altitude)}};
    \node[block, below=of raw, text width=52mm] (hyst) {zoom-change hysteresis\\{\footnotesize new zoom must hold ${\sim}2$\,s before it commits}};
    \node[block, below=of hyst, text width=52mm] (thresh) {zoom $<$ threshold?};
    \node[block, right=10mm of thresh, text width=34mm] (all) {drop tile filter,\\keep vehicle-exclude filter};
    \node[block, below=of thresh, text width=52mm] (same) {same 3$\times$3 tile set as loaded?};
    \node[block, below=of same, text width=52mm] (deb) {debounce: camera stable ${\sim}0.75$\,s};
    \node[block, below=of deb, text width=52mm] (refresh) {refresh only changed tiles\\{\footnotesize fetch new @ each geotile precision; unload gone; protect in-view/open as orphans}};
    \draw[flow] (raw) -- (hyst);
    \draw[flow] (hyst) -- (thresh);
    \draw[flow] (thresh) -- node[above,font=\footnotesize]{yes} (all);
    \draw[flow] (thresh) -- node[right,font=\footnotesize]{no} (same);
    \draw[flow] (same) -- node[right,font=\footnotesize]{no} (deb);
    \draw[flow] (deb) -- (refresh);
  \end{tikzpicture}
  \caption[The tile-refresh decision path]{The per-check tile-refresh decision path in \texttt{ADT4MOBGamemode}. Each guard -- hysteresis, threshold, no-op detection, debounce, changed-tiles-only refresh, orphan protection -- was added in response to an observed misbehaviour.}
  \label{fig:tile-flow}
\end{figure}

Making this stable in the face of a continuously moving camera required several guards, each added in response to an observed misbehaviour:

\begin{itemize}
    \item a periodic check (every ${\sim}0.5$\,s, rather than per-frame) detects when the camera has entered a new tile;
    \item a debounce delays the actual refresh until the camera has held still for ${\sim}0.75$\,s, so sweeping the view across several tiles issues one refresh, not one per tile crossed;
    \item a refresh only destroys and re-fetches the tiles that actually changed, keeping the ones still in view;
    \item a zoom-change hysteresis (Listing~\ref{lst:hysteresis}) requires a new zoom level to hold for ${\sim}2$\,s before it is committed, because camera-altitude jitter sitting exactly on a zoom boundary would otherwise oscillate the level, and every flip rebuilds the entire scene;
    \item entities that are still in view or have an open detail window when their tile unloads are kept alive as ``orphans'' and swept later once they are neither, rather than being destroyed out from under the user.
\end{itemize}

\begin{lstlisting}[style=cpp,float=tp,caption={[The zoom-change hysteresis guard]The zoom-change hysteresis in \texttt{ADT4MOBGamemode::CheckAndRefreshTiles}. \texttt{AltitudeToZoomLevel} rounds a logarithm, so an altitude on a zoom boundary flips on sub-metre jitter; each flip would run \texttt{DestroyAllActors}.},label={lst:hysteresis}]
int32 Zoom = RawZoom;
if (LoadedZoom >= 0 && RawZoom != LoadedZoom)
{
    if (RawZoom == CandidateZoom)
    {
        ZoomHoldTimer -= DeltaSeconds;
        if (ZoomHoldTimer > 0.f)
            Zoom = LoadedZoom;          // not held long enough yet -- stay put
    }
    else
    {
        CandidateZoom = RawZoom;
        ZoomHoldTimer = ZoomChangeHoldSeconds;
        Zoom = LoadedZoom;
    }
}
else { CandidateZoom = -1; }
\end{lstlisting}

Different data sources in the backend were also found to encode an entity's geotile at different precisions -- simulated vehicles at one zoom, imported static infrastructure at another -- so every tile query is issued at each known precision and the results merged, rather than assuming one encoding.

The live event subscription is scoped to match: when zoomed out past the tile-filtering threshold, individual simulated vehicles are sub-pixel and their high-churn update stream would dominate the WebSocket load for no visible benefit, so the subscription filter excludes that namespace entirely at that range.

\subsection{Serialising Concurrent Tile Requests}
\label{sec:impl-serialise}

An intermittent failure appeared during this work: \texttt{GetThingsByGeotileBounds} would occasionally log a transport-level failure -- not an HTTP error status, but the request failing to complete at all. It was traced to the refresh logic firing every changed tile's request concurrently in a tight loop, which the backend host handled unreliably under bursts. The fix at the time was to serialise the fetches: a queue held the pending tiles and started each tile's request only from the previous one's completion callback.

This was later removed. The residual flakiness was confirmed to be a server-side problem with the specific backend host, present before this code existed and independent of it, and once the backend was fixed the burst behaviour was no longer an issue and the current code fires the per-tile requests directly again. The episode is a useful case study for Section~\ref{sec:impl-discussion}: a decoupled architecture still has to respect the operational limits of whatever backend sits behind it, and a workaround written against one backend's transient behaviour is itself a form of coupling -- one that had to be carried until the backend changed and then deliberately unwound.

\subsection{Per-Entity Expiry for Simulated Fleets}
\label{sec:impl-ttl}

Ephemeral entities carry an expiry timestamp. For simulated vehicle fleets, the data source stamps the entire fleet with one shared timestamp, so honouring it literally destroys every vehicle in the same frame when the deadline passes, regardless of which ones are still being simulated. The generic path could not express ``expire these individually''. This is handled through the stateless extension point: the vehicle type's \texttt{UEntityTypeExtension} returns a per-entity staleness timeout (60\,s) instead, so an actor is removed only when no update for that specific entity has arrived within that window, with the timer reset on every message. Entities that genuinely stopped updating are removed, staggered naturally in time; entities still being simulated are not. The same extension also opts the type out of the per-message reflected-struct rebuild (\texttt{NeedsStructSyncOnLiveUpdate}), since a simulated vehicle's position is read straight from its raw JSON and the reflection walk -- a full traversal of the struct per message -- would be pure overhead at that update frequency. Both decisions are visible in Listing~\ref{lst:vehicleext} in the previous chapter.

\section{Adding an Entity Type: A Worked Example}
\label{sec:impl-worked-example}

To make the generalisation concrete, consider adding a hypothetical new infrastructure type -- say, an air-quality sensor with a \texttt{thingId} containing \texttt{air-quality} and attributes for the current readings. The full set of changes is:

\begin{enumerate}
    \item \textbf{Schema.} Add \texttt{AirQualityStruct.h} with \texttt{USTRUCT}s mirroring the sensor's JSON: a top-level struct with the latitude, longitude, and an \texttt{attributes} sub-struct holding the readings. \texttt{FJsonObjectConverter} handles deserialisation from that declaration alone.
    \item \textbf{Registration.} Add one line to \texttt{RegisterAllEntityTypes}: \texttt{Factory.RegisterType("air-quality", FAirQualityData::StaticStruct(), TEXT("Air Quality Sensor"), \ldots)}.
    \item \textbf{(Optional) presentation.} If the default grey badge and empty info panel are not wanted, add a \texttt{UAirQualityExtension} under \texttt{EntityDependencies/AirQuality/} overriding \texttt{GetBadgeColor}, \texttt{GetBadgeLabel}, and \texttt{GetDefaultInfoFields}.
\end{enumerate}

Nothing else changes. The factory now spawns an \texttt{ATempUIActor} for any matching \texttt{thingId}; the tile-streaming loop fetches it like any other entity; \texttt{UEntityUpdateDaemon} routes its live updates; the actor patches its JSON, updates its position through the coordinate seam, and refreshes the UI; selection, hover, and both camera modes work unmodified. There is no edit to the spawn pipeline, the update loop, the renderer, or the camera. This three-step (often two-step) process is the operational meaning of the generalisation claim in RQ1.

\section{Camera and Interaction}
\label{sec:impl-camera}

The camera is a single pawn, \texttt{AUnifiedPawn}, with two modes selected by an \texttt{ECameraMode} enum. In the RTS (overhead) mode it holds a fixed downward pitch with a spring arm whose length is the zoom, ranging from roughly $1.5$\,km to planetary distance, and re-aligns itself to the globe's local north and east every tick via \texttt{UpdateRtsAlignment}. In the free-fly mode it collapses the spring arm for full six-degree-of-freedom movement, staying upright relative to the globe surface by querying the local up-vector from the georeference in \texttt{UpdateFreeFlyAlignment}. Both modes' globe alignment is computed by projecting the ECEF north direction onto the local tangent plane.

The point worth making about the camera is what it does \textit{not} reference. It queries the globe -- the georeference -- for orientation, and it operates on actors and on the cursor for interaction. It has no knowledge of entity types, of the backend, or of the data model at all. Selection and hover are likewise generic: \texttt{USelectionManager}, a local-player subsystem, broadcasts hover and selection changes, and \texttt{ATempUIActor} responds by toggling custom depth-stencil values on all of its primitive components, with no per-type branching; the HUD widget binds to the same broadcasts to show or hide the entity detail window. This is precisely why the entire camera and interaction layer needed no changes while the entity side was being repeatedly restructured -- it was never coupled to the entity side in the first place.

\section{Special-Case Handling}
\label{sec:impl-special}

Fire-simulation entities are the one place genuinely infrastructure-specific behaviour exists in the viewer. A fire entity loads two separate multi-mesh models (a spread cone and a simulation body), fetches and parses GeoJSON perimeter geometry, and projects a precise terrain-exclusion polygon onto the globe so the burned area reads correctly against the terrain.

None of this could be expressed through a table entry or a stateless extension, because it involves live per-instance state: fetched GeoJSON, resolved model URLs, parsed vertex rings. It is therefore implemented as a \texttt{UEntityBehaviorComponent} subclass, \texttt{UFireBehaviorComponent} (Listing~\ref{lst:fire}), returned by the fire type's extension and attached generically by \texttt{ATempUIActor::Initialize} to each fire actor. The actor does not know what the component does; it just knows that some types supply one. The component overrides the two narrow points it needs -- it declares that it handles its own model loading, and it supplies exact exclusion polygons in place of the actor's default behaviour of hulling the visible mesh vertices.

\begin{lstlisting}[style=cpp,float=tp,caption={[The fire behaviour component]\texttt{UFireBehaviorComponent}. It is the only entity type that reaches the third rung of the entity model. The generic actor attaches it without knowing what it does.},label={lst:fire}]
UCLASS()
class DT4MOB_API UFireBehaviorComponent : public UEntityBehaviorComponent
{
    GENERATED_BODY()
public:
    virtual void OnEntityInitialized() override;   // load cone/sim GLBs, fetch GeoJSON
    virtual void OnEntityDataChanged() override;   // re-fetch on live perimeter update
    virtual bool HandlesOwnModelLoading() const override { return true; }
    virtual bool GetExclusionPolygons(
        TMap<FString, TArray<FVector2D>>& OutPolygons) const override;  // "Cone", "Simulation"
};
\end{lstlisting}

This section also carries a two-step history that illustrates the generalisation boundary being found in practice, sketched in Figure~\ref{fig:fire-fix}. The first refactoring separated the staged model loads for special-logic entities like fire out of the generic entity-load path, so that the generic path no longer had a fire-shaped branch in it. A later fix addressed a subtler coupling: the fire cone and the simulation exclusion polygon had been mutually exclusive on screen, not by design but as a side effect of sharing a single piece of state -- one \texttt{ACesiumCartographicPolygon} slot. The exclusion mechanism was changed so the behaviour component returns a \emph{map} of named polygons and \texttt{ATempUIActor} spawns one Cesium polygon per entry, so the cone and the simulation area can be shown at once. Both fixes pushed fire-specific behaviour further behind the same seam every other type uses, rather than special-casing it in the factory or the update loop.

\begin{figure}[t]
  \centering
  \begin{tikzpicture}[node distance=7mm, every node/.style={font=\small}]
    \node[block, text width=44mm] (b1) {\textbf{before}\\{\footnotesize one polygon slot shared}\\cone \emph{xor} simulation visible};
    \node[block, right=16mm of b1, text width=48mm] (a1) {\textbf{after}\\{\footnotesize \texttt{TMap<FString,TArray<FVector2D>>}}\\one Cesium polygon spawned per entry\\cone \emph{and} simulation visible};
    \draw[flow] (b1) -- (a1);
  \end{tikzpicture}
  \caption[The exclusion-polygon coupling fix]{Removing an accidental coupling: the cone and simulation areas were mutually exclusive only because they shared one polygon slot. Keying them by name in a map removed the constraint without any change to the generic actor beyond iterating the map.}
  \label{fig:fire-fix}
\end{figure}

\section{Discussion}
\label{sec:impl-discussion}

Some of these decisions generalised cleanly and some introduced complexity that is worth naming.

The registration table plus the two extension points held up well. Every entity type added after the pattern was in place fit one of its three rungs, and the generic spawn, update, and render paths did not need to change to accommodate any of them -- the worked example in Section~\ref{sec:impl-worked-example} is the normal case. The fire component -- the most demanding entity in the system -- fit the stateful-component rung without forcing an exception. The coordinate-conversion seam and the configuration-as-data change were similarly low-cost and paid off immediately, the latter by removing an entire class of ``which service parsed the config'' bug.

The decoupling from the backend is real but partial, as Section~\ref{sec:impl-replace-caveat} sets out: backend knowledge is confined to one layer, but that layer is a concrete implementation, not an interface with alternatives behind it, and no offline mock was built. Backend independence is therefore demonstrated by construction and by inspection -- one can point to every place the backend is named and observe that they are all in the service layer -- rather than by a working substitution.

The tile-streaming loop is where complexity accumulated. Each of its guards -- debounce, changed-tiles-only refresh, zoom hysteresis, orphan sweep, multi-precision geotile queries -- is individually justified by an observed failure, but together they make that loop the most intricate part of the codebase, and its correctness rests on empirically determined constants (hold times, thresholds) rather than on a clean invariant. The request-serialisation workaround, added and then removed, is the clearest example of the architecture being unable to fully insulate itself from backend behaviour: for a period, the viewer's fetch logic was shaped by one host's transient reliability rather than by its own needs.

This sets up the evaluation in Chapter~\ref{chap:evaluation}. The performance benchmarks target the entity spawning and rendering path -- exercised most heavily by high-cardinality types such as signs and simulated vehicles -- to establish how it scales with count and which rendering strategy such a type should use (RQ3); the per-type indirection itself (one map lookup for the extension, one for the schema, per spawn) is negligible by inspection and is not separately measured. The user testing targets the camera and interaction layer, since the question there is whether a layer built to be generic over the data still feels coherent to someone using it to explore a specific dataset (RQ4).

\chapter{Evaluation}
\label{chap:evaluation}

% TODO: draft, expand once benchmarks are finalized and user testing has run
\begin{introduction}
This chapter evaluates the viewer through performance benchmarks and user testing.
\end{introduction}

\section{Performance Benchmarks}
% TODO: sign model spawning benchmark harness and results (each metric
% averaged over 10 runs per supervisor feedback)

\section{User Testing}
% TODO: NOT YET RUN. Study design/materials are ready; write this section
% once the study has been conducted. Likely last section written.

\section{Discussion}
% TODO: synthesize benchmark + user testing results, relate back to the
% generalization/decoupling claims from Chapters~\ref{chap:architecture}
% and~\ref{chap:implementation}

\chapter{Conclusions}
\label{chap:conclusions}

\begin{introduction}
This chapter revisits the research questions in light of the work presented, summarises the contributions, honestly sets out the limitations of the current system, and describes the future work that follows most directly from it, including the viewer's intended release as open-source tooling for other researchers.
\end{introduction}

\section{Revisiting the Research Questions}
\label{sec:conc-rqs}

\paragraph{RQ1: a single pattern for arbitrary entity types.} This is answered affirmatively. Chapter~\ref{chap:architecture} presented an entity model built on one registration table and two escalating extension points, and Chapter~\ref{chap:implementation} showed that roughly thirty entity types, including static infrastructure, telematics equipment, simulated vehicles, and fire simulations, are supported through it, with all but one requiring no bespoke code in the generic spawn, update, or rendering paths. The worked example in Section~\ref{sec:impl-worked-example} showed that adding a further type is a two- or three-step change confined to a schema declaration, one table line, and an optional presentation class. The most demanding case, fire, fits the deepest extension point without forcing an exception.

\paragraph{RQ2: backend independence.} This is answered both by construction and by an actual interface. Backend knowledge is confined to the service layer described in Chapter~\ref{chap:architecture}, verifiable by inspection since no code above it names an HTTP verb, a URL, an authentication scheme, or the live-event envelope format; and, as Section~\ref{sec:impl-backend-interface} shows, that confined knowledge sits behind three small interfaces (for authentication and record access, for the live-connection protocol, and for decoding live-update messages) rather than only inside concrete classes. Pointing the viewer at a different backend is therefore writing one implementation of each interface and using it in place of the Ditto one, not rewriting the backend connection, the live-update stream, or the update router; describing that backend's entities is separate, ordinary work, the same work every new entity type already requires regardless of which backend supplies it (Section~\ref{sec:usecases}). No second implementation or offline mock has been written yet. Hence, the claim still rests on the existence and shape of the interfaces rather than a working substitution, but that remaining gap is a bounded implementation task, not an architectural one.

\paragraph{RQ3: scaling and rendering strategy.} This is answered. The ten-run sweep (Chapter~\ref{chap:evaluation}) rules out the naive \textsc{PerSign} strategy decisively: at 15,000 instances it issues 15--23\(\times\) more draw calls, spawns 37--38\(\times\) slower, and runs 50\% slower per frame than the two batched alternatives. Between those two, the result is narrower than the ``always atlas'' assumption this benchmark set out to test: \textsc{AtlasMaterial}'s much lower draw-call count makes it faster than \textsc{PerTexture} by less than 1\% of the frame time at the smaller counts, a statistically significant but practically negligible difference, and the two are not significantly different at 15,000 instances. At the same time, \textsc{PerTexture} wins decisively on memory (22\(\times\) less texture memory). The recommendation for this viewer's actual sign rendering is \textsc{PerTexture}, not the atlas approach the literature default would suggest.

\paragraph{RQ4: usability.} Answered positively, with qualifications. In a moderated study with ten participants from the intended user population (Section~\ref{sec:eval-user-testing}), the viewer scored a mean SUS of 81.1, between ``good'' and ``excellent'', and every participant completed the tasks covering camera control, entity selection, the detail panel, the geotechnical slope, the fire simulation, and navigation, with one performance-related failure on the ring-road task. The traffic task could not be run because the live traffic feed was unavailable, so that use case's usability is untested. The issues raised, chiefly keyboard-first navigation, poor entity visibility at low zoom, and missing feedback when placing an entity, belong to the camera and UI layers rather than to the decoupled design, and are addressable without architectural change.

\section{Summary of Contributions}
\label{sec:conc-contributions}

The four contributions stated in Chapter~\ref{chap:introduction} were delivered as follows. The first is an open-source digital twin viewer for mobility infrastructure, built on Unreal Engine and Cesium for Unreal, that renders roughly thirty entity types on a full-scale globe with live updates, with the entity model and backend-access layer isolated behind explicit seams. The second is the architectural pattern behind it: a service layer of small, single-purpose services that confines all backend contact, with the Ditto-specific pieces of that contact behind three small interfaces, and an entity model driven by a registration table with a stateless per-type extension point and a stateful per-instance component point, documented through the implementation decisions that realise it.

The third is a benchmark harness for the entity spawning and rendering path, comparing three mesh and material batching strategies for a high-cardinality type across a swept instance-count range, with each configuration averaged over ten runs. It rules out \textsc{PerSign} on draw calls, spawn time, and frame time, and recommends \textsc{PerTexture} over \textsc{AtlasMaterial} for this viewer's sign rendering on memory grounds, since the two are practically equivalent on frame time. The fourth is a user study with ten participants from the intended user population, evaluating the viewer through eight realistic tasks and the System Usability Scale.

\section{Limitations}
\label{sec:conc-limitations}

The main limitation is that backend decoupling has an interface but only one implementation. As set out under RQ2 and in Sections~\ref{sec:impl-backend-interface} and~\ref{sec:impl-replace-caveat}, the REST and authentication surface, the live-connection protocol, and the envelope format each sit behind their own interface. Still, exactly one implementation of each exists, no mock backend has been built, and the viewer has been validated against exactly one backend.

A second limitation is that the viewer inherits the backend's operational behaviour. Like any consumer application, it assumes the backend it points to behaves correctly; when the backend has connectivity or reliability problems of its own, the viewer is affected too, and no amount of code-level decoupling changes that.

\section{Future Work}
\label{sec:conc-future}

Future work falls into two groups: completing the architectural claims and making the viewer available beyond DT4MOB, and improving the application along the lines its users identified in the study in Section~\ref{sec:eval-user-testing}.

\subsection{Architecture and Release}

The most direct next step is to write a second implementation, or a mock, behind the existing interfaces. Section~\ref{sec:impl-backend-interface} describes the three backend interfaces, each with one Ditto implementation; writing a second implementation, or an offline mock that replays recorded traffic, would turn RQ2's architectural claim into a demonstrated one, and would let the viewer be developed and tested with no live backend at all. Beyond that, pointing the viewer at a different digital twin backend, or at a different Ditto deployment with a different entity schema, would be the strongest possible evidence for the decoupling.

A further step is to build an in-house, high-detail map of Portugal. Building the viewer's own map from high-detail geographic tiles, served as its own tile set, would give the project full control over the map's coverage, level of detail, and availability.

Finally, the viewer is intended for release as open-source software, so that other researchers and institutions working with digital twins of different infrastructure types can adopt it as a starting point rather than rebuilding a geospatial viewer from scratch. Such a release needs documentation of the extension points for external developers, a worked example repository adding a new entity type against a public dataset, a packaged build with the login flow pointed at a configurable backend, and a contribution guide. The architectural work in this thesis is what makes the release useful: a viewer that could only talk to one backend and show one project's entity types would not be worth adopting.

\subsection{Improvements from User Feedback}
\label{sec:conc-future-feedback}

The user study surfaced several concrete improvements. Almost all of them concern the camera and interface layers rather than the architecture, and most are inexpensive. They are discussed below by the part of the application they affect, in roughly the order of how often participants raised them.

The first concerns navigation. Keyboard-driven movement was the most frequent complaint, raised by roughly half the participants, who expected to pan and drag the globe with the mouse as in the map tools they use daily. The RTS camera should support click-and-drag panning (with the left or middle button) alongside the existing keys, so that the keyboard becomes optional rather than required. Participants also suggested rotation and vertical movement on familiar RTS-style keys, and one asked for camera movement to be clamped to the area of interest so the view cannot drift off the region with data. The FreeFly camera, in turn, was found too fast and visually heavy. It should offer an adjustable speed multiplier, reduce or make motion blur optional, and use Esc (rather than Tab) to release the cursor, which participants found more natural.

The second concerns finding entities at a distance. Participants found entities hard to see and to select when zoomed out, and hard to find without knowing in advance where they were. Three changes follow. Entity markers should grow with camera distance or have a minimum on-screen size, and use colour to distinguish types. Entities close together should collapse into a single cluster marker at far zoom levels, expanding as the camera approaches. A country-level overview should also mark where each digital twin's data is, so a user can jump to an area of interest instead of searching for it. To be consistent with the architecture, clustering and the overview should work on entity position alone, without per-type logic. A readout of the camera's current coordinates was also requested.

The third concerns feedback when placing entities and running simulations. Creating a fire simulation was the lowest-rated completed task, largely due to a lack of feedback. The placement flow should confirm that an entity was created, since one participant placed a second ignition point because they believed the first had failed. After \emph{Place Entity} is pressed, the cursor should change, or a prompt should appear, to show that a location is now expected; the entity type could also be chosen after pressing the button rather than before, which better matched participants' expectations. For the simulation itself, participants asked for a legend for the fire's colour scale and a progress indicator while it runs; one suggested a dedicated fire-simulation entry point rather than a generic placement tool.

The fourth concerns rendering quality. One participant's work machine could not run the viewer until a lighter build was made by changing engine rendering settings alone (disabling ray tracing and heavier anti-aliasing, among others), and some participants still saw slowdowns, notably after creating a fire simulation and when zooming in over the Porto region (Section~\ref{sec:eval-discussion}). Offering those settings as user-selectable quality presets, instead of fixing them at build time, would let one build serve both workstations and ordinary work machines.

Beyond these, one participant asked for dashboards that aggregate data across entities for analysis. Such dashboards already exist in Grafana, outside the viewer, and some entities link to them from their detail window; the request was to have them inside the viewer itself. Embedding them, or showing aggregated views alongside the map, would keep spatial inspection and analysis in one place.

Finally, the evaluation itself should be completed. The traffic task could not be run because the live feed was unavailable, so the traffic use case's usability is still untested; it should be re-run once the feed is back. A follow-up study after the changes above would also show whether they address the issues raised. It should record participants' hardware to separate interface problems from performance problems, and time each task, which this study did not.

\section{Closing Remarks}
\label{sec:conc-closing}

The central claim of this thesis is that a digital twin viewer for physical infrastructure need not be a closed, single-purpose application. By confining backend knowledge to a service layer and infrastructure-domain knowledge to a small set of opt-in extension points, the bulk of the viewer, meaning rendering, interaction, the camera, and coordinate placement, can be written once and left untouched as entity types are added and, in principle, as the backend changes. This is not free: the implementation in Chapter~\ref{chap:implementation} carries its own complexity, most visibly in the tile-streaming loop. But the result is a viewer whose cost of change is localised, and which is a plausible basis for reuse beyond the project it was built for.


%%%%%%%%%%%%%%%%%%%%%%%%%%%%%%%%%%%%%%%%%%%%%%%%%%%%%%%
% Print all used references
%%%%%%%%%%%%%%%%%%%%%%%%%%%%%%%%%%%%%%%%%%%%%%%%%%%%%%%

\begingroup
\renewcommand{\bibfont}{\footnotesize}
% References heading (English)
\setsecnumdepth{none}
\maxsecnumdepth{none}

\defbibheading{bibliography}[References]{
\chapter{References}
}
\vspace{2em}
\SingleSpacing
\setlength\bibitemsep{8pt}
\printbibliography[heading=bibliography]

\setsecnumdepth{subsubsection}
\maxsecnumdepth{subsubsection}
\endgroup


%%%%%%%%%%%%%%%%%%%%%%%%%%%%%%%%%%%%%%%%%%%%%%%%%%%%%%%
% Load appendix
%%%%%%%%%%%%%%%%%%%%%%%%%%%%%%%%%%%%%%%%%%%%%%%%%%%%%%%

\mainmatterWithoutReset
\appendix
\pagestyle{plain}
\chapter{User Test Materials}
\label{chap:appendix-user-test}

This appendix reproduces, in the original Portuguese, the materials used in the user study of Section~\ref{sec:eval-user-testing}: the invitation sent to participants, the task questionnaire, which includes the written introduction to the system that participants read before the tasks, and the final evaluation questionnaire. Links and personal contact details have been removed.

\section{Invitation}
\label{sec:appendix-invitation}

The invitation was sent to the DT4MOB project coordinator, who forwarded it to the participants from the partner companies and later sent a reminder.

\begin{otherlanguage}{portuguese}
\begin{quote}
\small
\textbf{Assunto:} DT4MOB -- Teste de usabilidade

Olá,

No âmbito do projeto DT4MOB, tenho realizado trabalho no desenvolvimento da interface (UI) do DT4MOB Twin Viewer, a aplicação que permite visualizar e interagir com a informação do Digital Twin sem requirir conhecimentos técnicos de programação.

Estamos agora numa fase de testes com utilizadores e gostaríamos de contar com a sua participação. O teste (15 a 30 mins) será online por video chamada e consiste em correr a aplicação e realizar um conjunto de tarefas simples associadas com os vários use cases da aplicação que consistem, em navegação, visualização de informação, e interação com o digital twin; para percebermos se a interface é intuitiva, se a informação apresentada é clara e útil, e para identificarmos problemas que só se revelam com uso real. Se estiver disponível para participar, agradecemos que se inscreva através do link: [ligação para marcação da sessão]

Muito obrigado desde já pelo seu tempo e contributo!

Links extras:
\begin{itemize}
    \item App Download: [ligação para download da aplicação]
    \item Forms do Teste -- Teste de Usabilidade: [ligação para o questionário de tarefas]
    \item Forms de avaliação -- Avaliação de Usabilidade: [ligação para o questionário de avaliação]
\end{itemize}

Com os melhores cumprimentos,\\
Bernardo Figueiredo
\end{quote}

\noindent Mensagem de encaminhamento do coordenador do projeto:

\begin{quote}
\small
Boa tarde,

Encaminho este e-mail do Bernardo Figueiredo, que vos convida a testar e avaliar a interface (UI) do DT4MOB. Esta vossa avaliação é muito importante para os trabalhos de mestrado do Bernardo, mas é ainda mais relevante para avaliarmos o estado de implementação do DT4MOB e identificarmos potenciais linhas de melhoria.

Obrigado,\\
{[coordenador do projeto DT4MOB]}
\end{quote}

\noindent Lembrete posterior do coordenador do projeto:

\begin{quote}
\small
Bom dia,

Reforço o pedido para que participem nos testes de usabilidade da interface do nosso projeto. É importante identificarmos pontos de melhoria para a nova iteração do trabalho.

Obrigado,\\
{[coordenador do projeto DT4MOB]}
\end{quote}
\end{otherlanguage}

% Form pages are cropped images of the printed forms (the browser's print
% header with the form URL, and the footer with page count and date, are
% cropped away and do not survive as hidden text).
\newcommand{\formpage}[1]{%
  \begin{center}
    \fbox{\includegraphics[height=0.8\textheight]{#1}}
  \end{center}
  \clearpage}

\clearpage
\section{Task Questionnaire}
\label{sec:appendix-task-form}

\formpage{appendix-task-form-1.png}
\formpage{appendix-task-form-2.png}
\formpage{appendix-task-form-3.png}
\formpage{appendix-task-form-4.png}
\formpage{appendix-task-form-5.png}
\formpage{appendix-task-form-6.png}
\formpage{appendix-task-form-7-anon.png}
\formpage{appendix-task-form-8-anon.png}

\section{Evaluation Questionnaire}
\label{sec:appendix-eval-form}

\formpage{appendix-eval-form-1.png}
\formpage{appendix-eval-form-2.png}
\formpage{appendix-eval-form-3.png}

\end{document}


